\section{Limitations of the Prototype}
\label{sec:limitations-of-prototype}

\minisec{Eingegrenzter Funktionsumfang}
Der Prototyp ist im Funktionsumfang bewusst begrenzt und als Minimum Viable Product primär auf
Hidden Profiles zugeschnitten, damit Fokus und anschließende Validierung zurechenbar bleiben
(vgl.\ \autoref{sec:hidden-profiles}).
Diese Eingrenzung setzt voraus, dass sich das Phänomen isoliert betrachten lässt. Diese Voraussetzung
ist vermutlich unterkomplex: In realen Meetings überlagert sich der Shared-Information-Bias mit Status- und
Hierarchieeffekten, Confirmation Bias und Groupthink sowie Zeitdruck, thematischer Abweichung und
unklarer Zielsetzung (vgl.\ \autoref{sec:key-challenges-in-meetings}).
Dieses Zusammenwirken deckt die vorliegende Version, wenn überhaupt nur implizit ab.
Zusätzlich sind mit dem Fortschritt allgemein verfügbarer KI-Systeme Menschen daran gewöhnt, frei mit
Sprachmodellen zu sprechen und Antworten auf beliebige Anliegen zu erwarten.
Diese Erwartungshaltung kann das enge Repertoire des Prototyps übersteigen. Teilnehmende
könnten Antwort- und Interaktionsmöglichkeiten voraussetzen, die das System nicht bereitstellt und somit das
vertrauen, auf Grund dieser wahrgenommenen Lücke, auch in bestehende, eigentlich nicht mangelhafte Funktionen des
Systems verlieren.

\minisec{Unimodale Wahrnehmung}
Das System arbeitet ausschließlich auf dem Transkript.
Transkriptionsfehler, etwa bei undeutlicher Aussprache oder schlechter Mikrofonqualität, gehen direkt
in den Kontext der LLM-Pipeline ein und können Gedanken und Interventionen verzerren.
Prosodie (Betonung, Aussrpache Tonfall) und visuelle Signale wie Mimik, Blick und sichtbare Aufmerksamkeit
oder Desinteresse fehlen.
Facilitatoren nutzen jedoch genau diese Signale, um zu lesen, wie eine Aussage gemeint ist und ob jemand
präsent ist (vgl.\ \autoref{ch:empirical-analysis-of-facilitating-strategies}).
Eine Sentimentanalyse auf dem Audiostream wurde nicht implementiert.
Das ist eine interpretative Lücke. Sie steht aber im Einklang mit den in den Interviews geäußerten
Bedenken gegenüber Kontrolle und Überwachung.

\minisec{Laufzeitmechanik}
Die Thought-Pipeline ist mehrstufig und probabilistisch. Interventionen können deshalb verzögert
eintreffen, und dieselbe Gesprächslage kann unterschiedlich ausfallen.
FIM-Scores werden einmal vergeben und nur deterministisch gealtert. Sie werden nicht semantisch neu
bewertet (vgl.\ \autoref{sec:dynamic-im}).

\minisec{Dazwischenreden}
Spricht eine teilnehmende Person in eine laufende Äußerung des Agenten hinein, bricht das System
diese nicht ab.
Der bereits ausgewählte Gedanke wird bis zum nächsten menschlichen Redeanteil zu Ende gesprochen.
Erst dieses neue Element stößt die Pipeline erneut an, sodass ein weiterer Gedanke gebildet und
geäußert werden kann.
Eine überlappende Wortmeldung verändert den laufenden Beitrag daher nicht und kann zu merkwürdigen bzw.
unnatürlichen Interaktionen führen, wenn Teilnehmende weitersprechen, während der AI-Facilitator spricht.

\minisec{Eigene Meetingplattform}
Das Meeting läuft in einer eigens gebauten Webanwendung und nicht in etablierten Meetingplattformen wie z.B.
Zoom oder Teams. Das sichert Kontrolle über Interaktion und UX. Jedoch sind Teilnehmende diese neue
Anwendung im Gegensatz zu den bekannten Tools nicht gewohnt.
Die Vergleichbarkeit mit der bekannten Nutzungssituation ist deshalb eingeschränkt.

\minisec{Modifizierte Turn-Allocation-Prediction}
Die Vorhersage, wer als Nächstes spricht, stützt sich auf Liu et al.\
\cite{liuProactiveConversationalAgents2025}.
Der Prompt musste jedoch von Alltagsdialogen auf den Facilitationskontext angepasst werden.
Deshalb kann nicht davon ausgegangen werden, dass sich die Ergebnisse aus dem Paper 1 zu 1 auf diesen Kontext
übertragen lassen.
Eine genauere statistische Validierung wäre empfehlens wert, war aber mangels eines geeignene
domänenspezifisches Korpus aus Facilitationssettings mit Turn-allocation Notation, nicht möglich.

\minisec{Annahme gleicher Redeanteile}
Der Participation Balancer setzt voraus, dass eine ausgeglichene Verteilung der Redezeit dem
Meeting zuträglich ist (vgl.\ \eqref{eq:fair-slot}).
Diese Voraussetzung gilt für das evaluierte Entscheidungsformat, in dem ungeteilte Informationen
nur durch Wortmeldungen in die Diskussion gelangen.
Sie gilt nicht für Formate, deren Zweck eine ungleiche Redezeit ist, etwa Präsentationen oder
längere Beiträge von Expertinnen und Experten.
Dort würden Curtailment und Invitation auslösen, obwohl die ungleiche Verteilung beabsichtigt ist.
